% TOP_PROBLEM: {"id": 1330, "title": "Hadwiger–Nelson Problem", "category_id": 3, "category": {"id": 3, "name": "graph_theory", "display_name": "Graph Theory", "description": "Problems involving graphs, networks, and their properties.", "slug": "graph-theory", "order_index": 3, "created_at": "2026-07-31T15:26:25.671Z"}, "status": "open"}
% FIRST_POSTED: null
% ENTRY_KIND: statement_only
% Imported from: batch06.tex; UnsolvedMath ID 1330
\documentclass[11pt]{article}
\usepackage[margin=25mm]{geometry}
\usepackage{fontspec}
\setmainfont{Latin Modern Roman}
\usepackage{amsmath,amssymb,mathtools}
\usepackage{unicode-math}
\setmathfont{Latin Modern Math}
\usepackage{xurl,hyperref}
\hypersetup{colorlinks=true,urlcolor=blue}
\setcounter{secnumdepth}{0}
\setlength{\parindent}{0pt}
\setlength{\parskip}{5pt}
\setlength{\emergencystretch}{3em}
\newcommand{\sref}[2]{\hyperlink{source:#1:#2}{[S#2]}}
\newcommand{\cataloguescope}[1]{\paragraph{Recorded target.}#1}
\begin{document}
% UnsolvedMath ID 1330; source catalogue code GRAPH-006
\setcounter{section}{1329}
\section{Hadwiger–Nelson Problem}\label{1330}
{\small\sffamily UnsolvedMath ID 1330 · Graph Theory}
\subsection{Definitions and mathematical statement}

\paragraph{Shared notation.}$\mathbb N=\{1,2,\ldots\}$, and $\mathbb Z,\mathbb Q,\mathbb R,\mathbb C$ denote the integers, rationals, reals and complex numbers. Any other conventions are specified locally.


Equip $\mathbb R^2$ with its Euclidean distance. The unit-distance graph of the plane has every point as a vertex, with an edge between distinct points $x,y$ exactly when $\|x-y\|_2=1$. A proper coloring assigns different colors to the endpoints of every edge.
\paragraph{Statement recorded in the supplied source.}
Determine
\[
 \chi(\mathbb R^2)=\min\bigl\{k\in\mathbb N:\ \exists c:\mathbb R^2\to\{1,\ldots,k\},\quad
 \|x-y\|_2=1\Longrightarrow c(x)\ne c(y)\bigr\}.
\]
No measurability or continuity of $c$ is required. \sref{1330}{2}
\subsection{Short English statement}
Find the fewest colors needed to color all points of the Euclidean plane so that points exactly one unit apart have different colors.
\subsection{Sources}
\begin{description}
\item[\hypertarget{source:1330:1}{S1}] ULAM.AI and the UnsolvedMath Contributors, \emph{UnsolvedMath: A Curated Collection of Open Mathematics Problems}, record 1330 (GRAPH-006). CC BY 4.0. \url{https://huggingface.co/datasets/ulamai/UnsolvedMath}.
\item[\hypertarget{source:1330:2}{S2}] Aubrey D. N. J. de Grey, \emph{The chromatic number of the plane is at least 5} (2018), arXiv:1804.02385v3, Section 1, p. 1. \url{https://arxiv.org/abs/1804.02385v3}.
\end{description}
\label{end:1330}
\end{document}
